\section{Self-Improving Agents：反馈环不等于自主进化}

Preference optimization 在训练阶段更新 policy；agent 则在环境中循环感知、计划、调用工具、行动并接收反馈。课程把 self-improvement 分成 refinement、reflection、ReAct 与 tree search。这里的 improvement 常是单任务或单会话的行为改善，不应自动改写成模型永久学会了新能力。

\subsection{Agent 的最小定义}

上一章的方法主要在一个请求内部增加推理或更新参数；agent 则把模型接入持续存在的 state 和 environment。定义 agent 时最容易犯的错误是只画一个循环箭头，却不写观察来自哪里、动作有什么权限、memory 是否可信、何时停止。下面的状态方程把这些对象显式化，后续 refinement 和 tree search 才有清楚的落点。

一个 agent 可抽象为
\[
s_{t+1}=F(s_t,o_t,m_t),\qquad a_t=\pi_{\theta}(s_t,\mathcal{T}),
\]
其中 $o_t$ 是观察，$m_t$ 是 memory，$\mathcal{T}$ 是工具集合，$a_t$ 是动作，环境再返回下一观察。基础模型只是 $\pi_{\theta}$ 的一部分；权限、状态、工具 schema、错误恢复和停止条件同样决定系统行为。


\lecturefigure{slide-085.jpg}{What is an AI Agent：环境、模型、目标、工具、记忆与动作}{CS25 V5 official deck, p.~85}

\teachervoice{00:34:20--00:35:20，Chelsea 将 agent 定义为能感知环境、决策并行动的系统，并列出 memory、tools 与 feedback。课堂提示：把一次模型调用包装成函数并不足以获得 agent reliability。}

\subsection{Refinement 与 reflection}

Refinement 让模型生成答案、找弱点、再修订；reflection 更强调从 failure 中抽取可复用经验并写入 memory。二者都可能陷入 confirmation loop：同一模型既作答又审查，错误假设可能被反复强化。


\lecturefigure{slide-086.jpg}{Self-improvement：refinement、reflection、探索与多路径推理}{CS25 V5 official deck, p.~86}


\lecturefigure{slide-087.jpg}{Refinement：生成、评估、反馈和迭代修订}{CS25 V5 official deck, p.~87}


\lecturefigure{slide-088.jpg}{Reflection：从失败中提取经验并更新后续策略}{CS25 V5 official deck, p.~88}

\begin{warningbox}{Self-critique 需要外部锚点}
没有 tool outcome、test、human check 或独立 verifier 时，``我检查过了'' 只是另一段生成。工程上应把 critique 转成可执行断言，再由环境验证；无法验证的部分必须保留 uncertainty。
\end{warningbox}

\subsection{ReAct 与 LATS：把 reasoning 接到 action/search}

ReAct 在 thought 与 action 之间交替，使模型根据 observation 更新计划。LATS 把单路径扩展为多条 action trajectory，并用 MCTS-like selection 聚合 feedback；它提高探索，却成倍增加工具调用、状态管理与错误分支成本。


\lecturefigure{slide-089.jpg}{ReAct：reasoning、action 与 observation 交替}{CS25 V5 official deck, p.~89}


\lecturefigure{slide-090.jpg}{LATS：多路径 agent search、反馈聚合与树搜索}{CS25 V5 official deck, p.~90}

\begin{warningbox}{Agent improvement 的可靠性协议}
每次 iteration 都应保存输入 state、模型计划、实际 action、tool output、verifier evidence 和 memory write；否则系统失败后无法区分规划错、工具错、权限错还是记忆污染。对 refinement，测修订后正确率与新引入错误；对 reflection，测记忆是否在相似未来任务中真正提升且不会传播错误；对 ReAct，测工具参数、observation parsing 和停止条件；对 LATS，测搜索预算增加带来的边际收益与 evaluator bias。生产环境还需要 action idempotency、transaction rollback、secret isolation 和人工升级路径。所谓 self-improving 只有在可审计 evidence 上持续减少失败，才是工程能力，而不是模型生成了更多自评文字。
\end{warningbox}

\begin{knowledgebox}{读图：四种 self-improvement 改写不同状态}
Refinement 改写当前 answer；reflection 把失败经验写入 memory；ReAct 用 observation 更新当前 trajectory；LATS 在 search tree 上维护多条候选及其价值。Slide 90 的树结构支持“探索更多路径”，但不证明搜索评分可靠。读图时先标记 state、action、feedback 和 persistent memory，再问错误能否被环境验证。若 feedback 仍由同一模型无锚点生成，循环可能只是把原错误写得更完整。
\end{knowledgebox}

\teachervoice{00:35:20--00:39:10，课堂依次区分 refinement、reflection、ReAct 与 LATS。实践经验：这些机制的差异在 feedback 被写到哪里——当前文本、长期 memory、环境 state 还是 search tree。}

\begin{lstlisting}[caption={带验证和权限边界的 agent loop}]
state = initialize(goal, permissions, budget)
while not done(state):
    proposal = model.plan(state)
    action = policy_gate(proposal, permissions)
    observation = execute_in_sandbox(action)
    evidence = verify_outcome(observation)
    state = update_memory(state, action, evidence)
    if repeated_failure(state): escalate_or_stop(state)
\end{lstlisting}

\subsection{本章小结}

Agent 能力来自模型与环境闭环。Refinement 改写答案，reflection 写入经验，ReAct 连接工具，LATS 增加多路径搜索；真正可靠的 self-improvement 依赖可验证 outcome、权限控制、预算和失败停止，而不是无限自我对话。

